\documentclass[11pt]{article}
\usepackage[T1]{fontenc}
\usepackage{lmodern,microtype}
\usepackage[a4paper,margin=26mm]{geometry}
\usepackage{amsmath,amssymb,amsthm,booktabs}
\usepackage[colorlinks=true,linkcolor=blue,urlcolor=blue,citecolor=blue]{hyperref}
\newtheorem{theorem}{Theorem}
\newtheorem{proposition}[theorem]{Proposition}
\newcommand{\F}{\mathbb F_2}
\newcommand{\Oc}{\mathcal O}
\DeclareMathOperator{\rad}{rad}
\title{Orthogonality cliques over $\F$:\\a radical quotient proof and a coloring separation}
\author{Wenlin Zhang\thanks{National University of Singapore; The Omega Institute.}
\and Haobo Ma\thanks{ChronoAI Pte Ltd; The Omega Institute.}}
\date{Working note, 28 September 2026}
\begin{document}
\maketitle
\begin{abstract}
We determine the clique number of the graph on all nonzero subspaces of
$\F^n$, adjacent when they are mutually orthogonal for the standard dot
form, answering Nikandish's Problem~4.2. Summing the radicals of a clique
and passing to the resulting orthogonal quotient gives the sharp upper
bound. We also give explicit coloring certificates in dimension six:
the graph on lines has a $14$-coloring, whereas the full subspace graph
has chromatic number between $15$ and $16$. Thus the two chromatic numbers
already differ in this dimension.
\end{abstract}

\section{The clique number}

Let $V=\F^n$ with $B(x,y)=\sum_i x_i y_i$. Following Nikandish~\cite{N},
let $\Oc_n^*$ have all nonzero subspaces of $V$ as vertices. Distinct
$U,W$ are adjacent if $B(U,W)=0$. Let $N(r)$ denote the number of nonzero
subspaces of $\F^r$, with $N(0)=0$.

\begin{theorem}\label{thm:clique}
For every integer $n\ge1$,
\[
 \omega(\Oc_n^*)=\max\{n,\ N(\lfloor n/2\rfloor)+(n\bmod2)\}.
\]
\end{theorem}

\begin{proof}
Take a clique $\mathcal C$ and set
\[
 R=\sum_{U\in\mathcal C}\rad(U),\qquad
 \rad(U)=U\cap U^\perp,\qquad r=\dim R.
\]
Each summand is orthogonal to its own member and, by the clique condition,
to every other member. Consequently
\[
 U\subseteq R^\perp,\qquad R\subseteq R^\perp,\qquad
 U\cap R=\rad(U)\quad(U\in\mathcal C).
\]
For the last equality, one inclusion follows from the definition of $R$;
the other follows from $R\subseteq U^\perp$. Nondegeneracy of $B$ gives
$2r\le n$ and $\dim(R^\perp/R)=n-2r$.

The form descends to $Q=R^\perp/R$. Write $\overline U$ for the image of
$U$ in $Q$. Its restricted form is nondegenerate: if $u\in U$ represents
a vector annihilating $\overline U$, then $u\in\rad(U)\subseteq R$, so
its image is zero. The spaces $\overline U$ are pairwise orthogonal.
Their nonzero members form an indexed direct sum. Indeed, in a relation
$\sum_U x_U=0$ with $x_U\in\overline U$, pairing with arbitrary vectors
of one $\overline W$ and using its nondegeneracy gives $x_W=0$.
Thus there are at most $n-2r$ nonzero images. This argument also applies
when a restricted or quotient form is alternating.

The members with zero image are precisely those contained in $R$, and
there are at most $N(r)$ of them. Therefore
\begin{equation}\label{eq:bound}
 |\mathcal C|\le N(r)+n-2r.
\end{equation}
For $r=0$ this is $n$. For every $s\ge1$, embedding $\F^s$ as a hyperplane
in $\F^{s+1}$ gives
\[
 N(s+1)\ge N(s)+2:
\]
an outside line and the whole ambient space are two distinct additional
subspaces. Hence $N(s)-2s$ is nondecreasing for $s\ge1$. If $r\ge1$,
putting $t=\lfloor n/2\rfloor$ in \eqref{eq:bound} gives
\[
 |\mathcal C|\le N(t)+n-2t=N(t)+(n\bmod2).
\]
These two cases prove the upper bound.

The $n$ coordinate lines attain the first lower bound. For the second,
let $t=\lfloor n/2\rfloor$ and take
\[
 T=\operatorname{span}\{e_1+e_2,e_3+e_4,\ldots,e_{2t-1}+e_{2t}\}.
\]
This is a totally isotropic $t$-space. All its nonzero subspaces are
pairwise orthogonal and give a clique of size $N(t)$. When $n$ is odd,
$\dim T^\perp=t+1$, so $T^\perp$ is a distinct additional vertex
orthogonal to every subspace of $T$. This includes $n=1$, where $T=0$.
The two constructions attain both terms of the maximum.
\end{proof}

Nikandish~\cite{N} gives the coordinate and isotropic lower constructions,
including the extra perpendicular-space vertex at $n=7$. The argument
above supplies the general upper bound. In particular, it treats
overlapping and degenerate clique members without replacing them by lines.

\section{A strict coloring separation in dimension six}

Let $\Gamma_n$ be the graph on nonzero vectors of $\F^n$, adjacent when
their dot product is zero. It identifies with the induced graph on lines
in $\Oc_n^*$. The equalities of chromatic numbers at $n=4,5$ in~\cite{N}
do not continue to $n=6$.

\begin{proposition}\label{prop:color}
There are explicit certificates for
\[
 \chi(\Gamma_6)\le14<15=\omega(\Oc_6^*)\le\chi(\Oc_6^*)\le16.
\]
\end{proposition}

\begin{proof}
Represent a vector by the integer whose six binary digits are its
coordinates. The following fourteen sets partition $\{1,\ldots,63\}$.
Every two distinct vectors within a row have dot product one, so each
row is an independent set in $\Gamma_6$.
\begin{center}
\begin{tabular}{rl@{\qquad}rl}
\toprule
Color&Vectors&Color&Vectors\\
\midrule
1&$2,14,50,62,63$&8&$11,22,34,37,56$\\
2&$8,28,41,43,61$&9&$4,5,52,54,60$\\
3&$6,12,42,53,59$&10&$16,18,21,24$\\
4&$19,38,46,55$&11&$1,9,23,39,51$\\
5&$32,35,44,47,48$&12&$15,30,40,49$\\
6&$7,31,33,36,58$&13&$25,27,29,57$\\
7&$13,17,20,26,45$&14&$3,10$\\
\bottomrule
\end{tabular}
\end{center}
This proves $\chi(\Gamma_6)\le14$. The totally isotropic space
$T=\operatorname{span}\{e_1+e_2,e_3+e_4,e_5+e_6\}$ has seven lines,
seven planes and one three-dimensional subspace. These fifteen vertices
form the clique already exhibited in~\cite{N}; Theorem~\ref{thm:clique}
also gives its optimality.

For the upper bound on the full graph, the accompanying file
\path{coloring-certificate.json} assigns one of sixteen colors to each
nonzero subspace, specified by a basis. The independent checker enumerates
the span of each supplied basis and verifies independence and uniqueness.
The numbers of distinct subspaces in dimensions $1$ through $6$ are
\[
 63,\quad651,\quad1395,\quad651,\quad63,\quad1.
\]
These agree with the Gaussian subspace counts, so the list contains all
$2824$ nonzero subspaces. For each pair of distinct subspaces assigned
the same color, the checker finds a pair of basis vectors with dot
product one. Bilinearity therefore excludes orthogonality of the two
subspaces. This verifies a proper $16$-coloring of the complete graph.
\end{proof}

The exact values of $\chi(\Gamma_6)$ and $\chi(\Oc_6^*)$ are not inferred
from these colorings. The immediate remaining question for the full graph
is whether a $15$-coloring exists. A search can assign distinct fixed colors
to the fifteen subspaces of $T$, eliminating color-permutation symmetry,
and ask whether the precoloring extends to all remaining subspaces.
If no extension exists, a checkable obstruction would establish the value
$16$; an extension would establish $15$.

\section{Proof files and reproduction}

The complete Lean proof of Theorem~\ref{thm:clique} is supplied as
\path{NikandishClique.lean}. The public
\href{https://github.com/the-omega-institute/trureturing/blob/063bb60d1df414c93252023dbc5a5ae659566e58/D5/S3/Combinatorics/Orthogonality/NikandishClique.lean}{pinned source}
is in the trureturing repository at revision
\begin{center}\small\ttfamily
063bb60d1df414c93252023dbc5a5ae659566e58.
\end{center}
Its imports use
Lean~4.33.0 and the pinned Mathlib dependencies. The theorem is
\path{D5.S3.Combinatorics.Orthogonality.NikandishClique.result}.
The published formal verification records the standard axiom closure
\texttt{propext}, \texttt{Classical.choice}, and \texttt{Quot.sound}.

The new colorings are finite certificates checked with exact integer
arithmetic; they are not yet Lean formalized. From the accompanying package,
run \texttt{python3 check\_certificate.py}. To regenerate the certificates,
first run \texttt{python3 coloring\_certificate.py}. The checker does not
import the search program. AI tools assisted the exposition, code development
and computational exploration.

\begin{thebibliography}{1}
\raggedright
\bibitem{N} R. Nikandish, \emph{Annihilating-Ideal Graphs and Orthogonality
Graphs over $\mathbb F_2$}, arXiv:2609.22769v1 (2026).
\end{thebibliography}
\end{document}
